\documentclass[11pt]{article}
\newcommand{\videoid}{Ts7i-Am3nf8}
% Preámbulo compartido de las notas del curso de topología algebraica.
% Cada nota empieza con \documentclass[11pt]{article}, define \videoid{...}
% (el identificador del video de YouTube) y luego hace % Preámbulo compartido de las notas del curso de topología algebraica.
% Cada nota empieza con \documentclass[11pt]{article}, define \videoid{...}
% (el identificador del video de YouTube) y luego hace % Preámbulo compartido de las notas del curso de topología algebraica.
% Cada nota empieza con \documentclass[11pt]{article}, define \videoid{...}
% (el identificador del video de YouTube) y luego hace \input{preambulo}.
% Compilar con LuaLaTeX:  lualatex nombre-de-la-nota.tex
\usepackage[spanish,provide=*]{babel}
\usepackage[T1]{fontenc}
\usepackage{lmodern}
\usepackage[a4paper,margin=2.6cm]{geometry}
\usepackage{amsmath,amssymb,amsthm}
\usepackage{xcolor}
\usepackage{tikz}
\usetikzlibrary{arrows.meta,babel}
\usepackage[colorlinks=true,urlcolor=blue!50!black,linkcolor=blue!50!black]{hyperref}

\hyphenpenalty=10000 \exhyphenpenalty=10000 \emergencystretch=3em

% Colores de las figuras (los mismos que usa la página)
\definecolor{extremo}{RGB}{47,107,79}
\definecolor{azul}{RGB}{44,74,154}

% Definiciones numeradas por unidad
\theoremstyle{definition}
\newtheorem{defn}{Definición}
% Numeración por unidad: \unidad{u}{n} (después de \input{preambulo}) hace que las
% definiciones de la nota se llamen u.n, u.(n+1), ...; n es el primer número
% libre de esa unidad (ver la tabla de CONVENCIONES.md).
\newcommand{\unidad}[2]{\renewcommand{\thedefn}{#1.\arabic{defn}}\setcounter{defn}{\numexpr#2-1\relax}}

% Observaciones y erratas del video. Usar SOLO si hay algo que señalar:
%   \begin{erratas}
%     \errata{203}{3:23}{Se dice ..., pero debe decir ...}
%   \end{erratas}
% El primer argumento son los segundos (enlaza al minuto exacto del video),
% el segundo es la marca visible, el tercero es el texto.
\newenvironment{erratas}{\section*{Observaciones y erratas del video}\begin{itemize}}{\end{itemize}}
\newcommand{\errata}[3]{%
  \edef\errurl{https://www.youtube.com/watch?v=\videoid\string&t=#1}%
  \item[\textbf{\expandafter\href\expandafter{\errurl}{#2}}] #3}
.
% Compilar con LuaLaTeX:  lualatex nombre-de-la-nota.tex
\usepackage[spanish,provide=*]{babel}
\usepackage[T1]{fontenc}
\usepackage{lmodern}
\usepackage[a4paper,margin=2.6cm]{geometry}
\usepackage{amsmath,amssymb,amsthm}
\usepackage{xcolor}
\usepackage{tikz}
\usetikzlibrary{arrows.meta,babel}
\usepackage[colorlinks=true,urlcolor=blue!50!black,linkcolor=blue!50!black]{hyperref}

\hyphenpenalty=10000 \exhyphenpenalty=10000 \emergencystretch=3em

% Colores de las figuras (los mismos que usa la página)
\definecolor{extremo}{RGB}{47,107,79}
\definecolor{azul}{RGB}{44,74,154}

% Definiciones numeradas por unidad
\theoremstyle{definition}
\newtheorem{defn}{Definición}
% Numeración por unidad: \unidad{u}{n} (después de % Preámbulo compartido de las notas del curso de topología algebraica.
% Cada nota empieza con \documentclass[11pt]{article}, define \videoid{...}
% (el identificador del video de YouTube) y luego hace \input{preambulo}.
% Compilar con LuaLaTeX:  lualatex nombre-de-la-nota.tex
\usepackage[spanish,provide=*]{babel}
\usepackage[T1]{fontenc}
\usepackage{lmodern}
\usepackage[a4paper,margin=2.6cm]{geometry}
\usepackage{amsmath,amssymb,amsthm}
\usepackage{xcolor}
\usepackage{tikz}
\usetikzlibrary{arrows.meta,babel}
\usepackage[colorlinks=true,urlcolor=blue!50!black,linkcolor=blue!50!black]{hyperref}

\hyphenpenalty=10000 \exhyphenpenalty=10000 \emergencystretch=3em

% Colores de las figuras (los mismos que usa la página)
\definecolor{extremo}{RGB}{47,107,79}
\definecolor{azul}{RGB}{44,74,154}

% Definiciones numeradas por unidad
\theoremstyle{definition}
\newtheorem{defn}{Definición}
% Numeración por unidad: \unidad{u}{n} (después de \input{preambulo}) hace que las
% definiciones de la nota se llamen u.n, u.(n+1), ...; n es el primer número
% libre de esa unidad (ver la tabla de CONVENCIONES.md).
\newcommand{\unidad}[2]{\renewcommand{\thedefn}{#1.\arabic{defn}}\setcounter{defn}{\numexpr#2-1\relax}}

% Observaciones y erratas del video. Usar SOLO si hay algo que señalar:
%   \begin{erratas}
%     \errata{203}{3:23}{Se dice ..., pero debe decir ...}
%   \end{erratas}
% El primer argumento son los segundos (enlaza al minuto exacto del video),
% el segundo es la marca visible, el tercero es el texto.
\newenvironment{erratas}{\section*{Observaciones y erratas del video}\begin{itemize}}{\end{itemize}}
\newcommand{\errata}[3]{%
  \edef\errurl{https://www.youtube.com/watch?v=\videoid\string&t=#1}%
  \item[\textbf{\expandafter\href\expandafter{\errurl}{#2}}] #3}
) hace que las
% definiciones de la nota se llamen u.n, u.(n+1), ...; n es el primer número
% libre de esa unidad (ver la tabla de CONVENCIONES.md).
\newcommand{\unidad}[2]{\renewcommand{\thedefn}{#1.\arabic{defn}}\setcounter{defn}{\numexpr#2-1\relax}}

% Observaciones y erratas del video. Usar SOLO si hay algo que señalar:
%   \begin{erratas}
%     \errata{203}{3:23}{Se dice ..., pero debe decir ...}
%   \end{erratas}
% El primer argumento son los segundos (enlaza al minuto exacto del video),
% el segundo es la marca visible, el tercero es el texto.
\newenvironment{erratas}{\section*{Observaciones y erratas del video}\begin{itemize}}{\end{itemize}}
\newcommand{\errata}[3]{%
  \edef\errurl{https://www.youtube.com/watch?v=\videoid\string&t=#1}%
  \item[\textbf{\expandafter\href\expandafter{\errurl}{#2}}] #3}
.
% Compilar con LuaLaTeX:  lualatex nombre-de-la-nota.tex
\usepackage[spanish,provide=*]{babel}
\usepackage[T1]{fontenc}
\usepackage{lmodern}
\usepackage[a4paper,margin=2.6cm]{geometry}
\usepackage{amsmath,amssymb,amsthm}
\usepackage{xcolor}
\usepackage{tikz}
\usetikzlibrary{arrows.meta,babel}
\usepackage[colorlinks=true,urlcolor=blue!50!black,linkcolor=blue!50!black]{hyperref}

\hyphenpenalty=10000 \exhyphenpenalty=10000 \emergencystretch=3em

% Colores de las figuras (los mismos que usa la página)
\definecolor{extremo}{RGB}{47,107,79}
\definecolor{azul}{RGB}{44,74,154}

% Definiciones numeradas por unidad
\theoremstyle{definition}
\newtheorem{defn}{Definición}
% Numeración por unidad: \unidad{u}{n} (después de % Preámbulo compartido de las notas del curso de topología algebraica.
% Cada nota empieza con \documentclass[11pt]{article}, define \videoid{...}
% (el identificador del video de YouTube) y luego hace % Preámbulo compartido de las notas del curso de topología algebraica.
% Cada nota empieza con \documentclass[11pt]{article}, define \videoid{...}
% (el identificador del video de YouTube) y luego hace \input{preambulo}.
% Compilar con LuaLaTeX:  lualatex nombre-de-la-nota.tex
\usepackage[spanish,provide=*]{babel}
\usepackage[T1]{fontenc}
\usepackage{lmodern}
\usepackage[a4paper,margin=2.6cm]{geometry}
\usepackage{amsmath,amssymb,amsthm}
\usepackage{xcolor}
\usepackage{tikz}
\usetikzlibrary{arrows.meta,babel}
\usepackage[colorlinks=true,urlcolor=blue!50!black,linkcolor=blue!50!black]{hyperref}

\hyphenpenalty=10000 \exhyphenpenalty=10000 \emergencystretch=3em

% Colores de las figuras (los mismos que usa la página)
\definecolor{extremo}{RGB}{47,107,79}
\definecolor{azul}{RGB}{44,74,154}

% Definiciones numeradas por unidad
\theoremstyle{definition}
\newtheorem{defn}{Definición}
% Numeración por unidad: \unidad{u}{n} (después de \input{preambulo}) hace que las
% definiciones de la nota se llamen u.n, u.(n+1), ...; n es el primer número
% libre de esa unidad (ver la tabla de CONVENCIONES.md).
\newcommand{\unidad}[2]{\renewcommand{\thedefn}{#1.\arabic{defn}}\setcounter{defn}{\numexpr#2-1\relax}}

% Observaciones y erratas del video. Usar SOLO si hay algo que señalar:
%   \begin{erratas}
%     \errata{203}{3:23}{Se dice ..., pero debe decir ...}
%   \end{erratas}
% El primer argumento son los segundos (enlaza al minuto exacto del video),
% el segundo es la marca visible, el tercero es el texto.
\newenvironment{erratas}{\section*{Observaciones y erratas del video}\begin{itemize}}{\end{itemize}}
\newcommand{\errata}[3]{%
  \edef\errurl{https://www.youtube.com/watch?v=\videoid\string&t=#1}%
  \item[\textbf{\expandafter\href\expandafter{\errurl}{#2}}] #3}
.
% Compilar con LuaLaTeX:  lualatex nombre-de-la-nota.tex
\usepackage[spanish,provide=*]{babel}
\usepackage[T1]{fontenc}
\usepackage{lmodern}
\usepackage[a4paper,margin=2.6cm]{geometry}
\usepackage{amsmath,amssymb,amsthm}
\usepackage{xcolor}
\usepackage{tikz}
\usetikzlibrary{arrows.meta,babel}
\usepackage[colorlinks=true,urlcolor=blue!50!black,linkcolor=blue!50!black]{hyperref}

\hyphenpenalty=10000 \exhyphenpenalty=10000 \emergencystretch=3em

% Colores de las figuras (los mismos que usa la página)
\definecolor{extremo}{RGB}{47,107,79}
\definecolor{azul}{RGB}{44,74,154}

% Definiciones numeradas por unidad
\theoremstyle{definition}
\newtheorem{defn}{Definición}
% Numeración por unidad: \unidad{u}{n} (después de % Preámbulo compartido de las notas del curso de topología algebraica.
% Cada nota empieza con \documentclass[11pt]{article}, define \videoid{...}
% (el identificador del video de YouTube) y luego hace \input{preambulo}.
% Compilar con LuaLaTeX:  lualatex nombre-de-la-nota.tex
\usepackage[spanish,provide=*]{babel}
\usepackage[T1]{fontenc}
\usepackage{lmodern}
\usepackage[a4paper,margin=2.6cm]{geometry}
\usepackage{amsmath,amssymb,amsthm}
\usepackage{xcolor}
\usepackage{tikz}
\usetikzlibrary{arrows.meta,babel}
\usepackage[colorlinks=true,urlcolor=blue!50!black,linkcolor=blue!50!black]{hyperref}

\hyphenpenalty=10000 \exhyphenpenalty=10000 \emergencystretch=3em

% Colores de las figuras (los mismos que usa la página)
\definecolor{extremo}{RGB}{47,107,79}
\definecolor{azul}{RGB}{44,74,154}

% Definiciones numeradas por unidad
\theoremstyle{definition}
\newtheorem{defn}{Definición}
% Numeración por unidad: \unidad{u}{n} (después de \input{preambulo}) hace que las
% definiciones de la nota se llamen u.n, u.(n+1), ...; n es el primer número
% libre de esa unidad (ver la tabla de CONVENCIONES.md).
\newcommand{\unidad}[2]{\renewcommand{\thedefn}{#1.\arabic{defn}}\setcounter{defn}{\numexpr#2-1\relax}}

% Observaciones y erratas del video. Usar SOLO si hay algo que señalar:
%   \begin{erratas}
%     \errata{203}{3:23}{Se dice ..., pero debe decir ...}
%   \end{erratas}
% El primer argumento son los segundos (enlaza al minuto exacto del video),
% el segundo es la marca visible, el tercero es el texto.
\newenvironment{erratas}{\section*{Observaciones y erratas del video}\begin{itemize}}{\end{itemize}}
\newcommand{\errata}[3]{%
  \edef\errurl{https://www.youtube.com/watch?v=\videoid\string&t=#1}%
  \item[\textbf{\expandafter\href\expandafter{\errurl}{#2}}] #3}
) hace que las
% definiciones de la nota se llamen u.n, u.(n+1), ...; n es el primer número
% libre de esa unidad (ver la tabla de CONVENCIONES.md).
\newcommand{\unidad}[2]{\renewcommand{\thedefn}{#1.\arabic{defn}}\setcounter{defn}{\numexpr#2-1\relax}}

% Observaciones y erratas del video. Usar SOLO si hay algo que señalar:
%   \begin{erratas}
%     \errata{203}{3:23}{Se dice ..., pero debe decir ...}
%   \end{erratas}
% El primer argumento son los segundos (enlaza al minuto exacto del video),
% el segundo es la marca visible, el tercero es el texto.
\newenvironment{erratas}{\section*{Observaciones y erratas del video}\begin{itemize}}{\end{itemize}}
\newcommand{\errata}[3]{%
  \edef\errurl{https://www.youtube.com/watch?v=\videoid\string&t=#1}%
  \item[\textbf{\expandafter\href\expandafter{\errurl}{#2}}] #3}
) hace que las
% definiciones de la nota se llamen u.n, u.(n+1), ...; n es el primer número
% libre de esa unidad (ver la tabla de CONVENCIONES.md).
\newcommand{\unidad}[2]{\renewcommand{\thedefn}{#1.\arabic{defn}}\setcounter{defn}{\numexpr#2-1\relax}}

% Observaciones y erratas del video. Usar SOLO si hay algo que señalar:
%   \begin{erratas}
%     \errata{203}{3:23}{Se dice ..., pero debe decir ...}
%   \end{erratas}
% El primer argumento son los segundos (enlaza al minuto exacto del video),
% el segundo es la marca visible, el tercero es el texto.
\newenvironment{erratas}{\section*{Observaciones y erratas del video}\begin{itemize}}{\end{itemize}}
\newcommand{\errata}[3]{%
  \edef\errurl{https://www.youtube.com/watch?v=\videoid\string&t=#1}%
  \item[\textbf{\expandafter\href\expandafter{\errurl}{#2}}] #3}

\unidad{1}{29} % unidad 1; primer enunciado: 1.29

\title{Homotopías entre funciones}
\author{Luis Jorge Sánchez Saldaña\\[2pt]\small Facultad de Ciencias, UNAM}
\date{}

\begin{document}
\maketitle

\noindent\small Notas del video
\href{https://www.youtube.com/watch?v=Ts7i-Am3nf8}{\guillemotleft{}Homotopías entre funciones\guillemotright{}}
(curso de Introducción a la Topología Algebraica). Se define la homotopía entre funciones continuas y la homotopía que preserva el punto base, y se demuestra que funciones homotópicas (preservando el punto base) inducen el mismo homomorfismo en grupos fundamentales.
\normalsize

\section*{Homotopías entre funciones}

En este video se define la homotopía entre funciones, y también un tipo particular de homotopía que será útil al estudiar las funciones inducidas entre grupos fundamentales. Todas las funciones son continuas.

\begin{defn}[Homotopía entre funciones]\label{def:homotopia-funciones}
Sean $g,f\colon X\to Y$ dos funciones (continuas) con el mismo dominio y el mismo codominio. Se dice que $g$ es \emph{homotópica} a $f$, y se escribe $g\simeq f$, si existe una función continua $H\colon X\times I\to Y$, $(x,t)\mapsto H(x,t)$, tal que
\[
  H(x,0)=g(x),\qquad H(x,1)=f(x)\qquad\text{para todo }x\in X.
\]
\end{defn}

Aquí $I=[0,1]$ es el intervalo unitario y el segundo parámetro se piensa como el tiempo: en tiempo $0$ la función $H$ recupera $g$ y en tiempo $1$ recupera $f$, de modo que $g$ se \emph{deforma} en $f$ a través del tiempo. Esto generaliza la homotopía entre caminos: cuando $X=I$, se tiene justamente una homotopía entre caminos.

\begin{center}
\begin{tikzpicture}[x=0.022cm,y=-0.022cm,line cap=round]
  \draw[gray] (80,60) -- (80,200); \draw[gray] (240,60) -- (240,200);
  \draw[thick,black] (80,60) -- (240,60);
  \draw[gray] (80,95) -- (240,95);
  \draw[gray] (80,130) -- (240,130);
  \draw[gray] (80,165) -- (240,165);
  \draw[thick,azul] (80,200) -- (240,200);
  \draw[gray] (470,138) ellipse (105 and 92);
  \draw[thick] (385,90) .. controls (420,70) and (470,115) .. (545,90);
  \draw[gray] (385,115) .. controls (425,100) and (475,140) .. (545,118);
  \draw[gray] (385,138) .. controls (425,125) and (475,165) .. (545,142);
  \draw[gray] (385,160) .. controls (425,150) and (475,188) .. (545,166);
  \draw[thick,azul] (385,185) .. controls (425,175) and (475,210) .. (545,188);
  \draw[->,gray] (262,130) -- (340,130); \node at (301,112) {\small $H$};
  \node[anchor=east] at (72,62) {\small $X\times\{0\}$}; \node[anchor=east,azul] at (72,200) {\small $X\times\{1\}$};
  \node at (160,35) {$X\times I$}; \node at (470,22) {$Y$};
  \node at (565,75) {\small $g$}; \node[azul] at (565,198) {\small $f$};
\end{tikzpicture}
\end{center}

\noindent Figura: cada una de las líneas horizontales de $X\times I$ es una copia de $X$; la de arriba, $X\times\{0\}$, se manda con $H$ a la imagen de $g$, la de abajo, $X\times\{1\}$, a la de $f$, y las intermedias (en gris) a las etapas de la deformación.

\section*{Homotopías que preservan el punto base}

Ahora se define una variante de la definición anterior, muy útil para trabajar con homomorfismos inducidos en grupos fundamentales. Además de $X$ y $Y$ se consideran subespacios $A\subseteq X$ y $B\subseteq Y$. Primero, un poco de notación: si $f\colon X\to Y$ es una función tal que $f(A)\subseteq B$, se escribe
\[
  f\colon (X,A)\to(Y,B).
\]
De ahora en adelante, si se escribe esto, quiere decir que es una función de $X$ en $Y$ que manda $A$ en $B$. En particular, si $A=\{x_0\}$ y $B=\{y_0\}$ constan de un punto (los puntos base que se venían tomando), la función se escribe sin las llaves, $f\colon (X,x_0)\to(Y,y_0)$, y lo único que quiere decir es que $f$ manda el punto base $x_0$ en el punto base $y_0$. Como se vio en videos anteriores, estas funciones inducen homomorfismos entre los grupos fundamentales basados en $x_0$ y en $y_0$ (Teorema~1.25).

\begin{defn}[Homotopía entre funciones de pares]\label{def:homotopia-pares}
Sean $g,f\colon (X,A)\to(Y,B)$. Una \emph{homotopía entre $g$ y $f$ como funciones de pares} es una función continua $H\colon X\times I\to Y$ tal que, para cada $t\in I$,
\[
  H(\,\cdot\,,t)\colon (X,A)\to(Y,B),
\]
y además $H(x,0)=g(x)$ y $H(x,1)=f(x)$ para todo $x\in X$.

En particular, una homotopía entre $g,f\colon (X,x_0)\to(Y,y_0)$ es una función $H\colon X\times I\to Y$ como en la Definición~\ref{def:homotopia-funciones} tal que
\[
  H(x_0,t)=y_0\qquad\text{para todo }t\in I.
\]
Estas se llaman \emph{homotopías que preservan el punto base}.
\end{defn}

Es decir, es una deformación de $g$ en $f$ de manera que todas las etapas de la deformación siguen siendo funciones entre las parejas $(X,A)\to(Y,B)$; en el caso del punto base, $x_0$ se manda siempre en $y_0$ en todo momento. Cuando se dice que dos funciones $g,f\colon (X,x_0)\to(Y,y_0)$ son homotópicas, se entiende que lo son mediante una homotopía que preserva el punto base.

\section*{El resultado: homotopía e inducidos}

La gracia de estas homotopías es el siguiente resultado: funciones homotópicas que preservan el punto base inducen el mismo homomorfismo en grupos fundamentales.

\begin{prop}\label{prop:homotopia-inducido}
Si $g,f\colon (X,x_0)\to(Y,y_0)$ son homotópicas (mediante una homotopía que preserva el punto base), entonces
\[
  g_*=f_*\colon\pi_1(X,x_0)\to\pi_1(Y,y_0).
\]
\end{prop}

\begin{proof}
Tanto $g_*$ como $f_*$ son funciones con el mismo dominio y el mismo codominio, de modo que solo hay que verificar que tienen la misma imagen en cada elemento. Sea $[\gamma]\in\pi_1(X,x_0)$, con $\gamma$ un lazo en $X$ basado en $x_0$. Por definición, $g_*[\gamma]=[g\circ\gamma]$ y $f_*[\gamma]=[f\circ\gamma]$. Lo que se quiere demostrar es que los lazos $g\circ\gamma$ y $f\circ\gamma$, que viven en $Y$ y están basados en $y_0$, son homotópicos relativos a $\{0,1\}$.

Sea $H\colon X\times I\to Y$ una homotopía que preserva el punto base entre $g$ y $f$, que existe por hipótesis. Se considera la función $\gamma\times\mathrm{id}_I\colon I\times I\to X\times I$ y se compone con $H$:
\[
  H\circ(\gamma\times\mathrm{id}_I)\colon I\times I\to Y,\qquad (s,t)\mapsto H(\gamma(s),t).
\]
Esta composición es justamente la homotopía relativa a los extremos entre $g\circ\gamma$ y $f\circ\gamma$ que se buscaba. Verificarlo es un ejercicio sencillo: en $t=0$ y $t=1$ se obtienen $g\circ\gamma$ y $f\circ\gamma$, y para ver que es relativa a $\{0,1\}$ se usa que la homotopía $H$ siempre manda el punto base en el punto base, es decir, $H(\gamma(0),t)=H(\gamma(1),t)=H(x_0,t)=y_0$ para todo $t$.
\end{proof}

Lo que debe pensarse de esta proposición es que funciones homotópicas mediante homotopías que preservan el punto base inducen el mismo homomorfismo en grupos fundamentales.

\begin{erratas}
  \errata{336}{5:36}{En la definición de homotopía entre funciones de pares se escribe $H(x,0)=f(x)$ y $H(x,1)=g(x)$, al revés que en la definición anterior ($H(x,0)=g(x)$ y $H(x,1)=f(x)$). En estas notas se usa el mismo orden en ambas definiciones.}
\end{erratas}

\end{document}
